\section{Cote 136: what an escaping-prime obstruction cannot be}
\label{sec:136linearprimetail}

\paragraph{Archive, inputs and scope.}
We retain the precise category $\mathcal A_{\mathbb Z}$ of
AQ-136-COMPARISON: physical pp.~117 and 119, archive footers 116
and 118, \texttt{ARCHIVE\_136\_scan117.tex} and
\texttt{ARCHIVE\_136\_scan119.tex}. The question of components and
the all-arrow groupoid is physical p.~121, footer 120,
\texttt{ARCHIVE\_136\_scan121.tex}.
The following result extends the already supplied
AQ-136-DERIVED-COMPATIBILITY. It does not give a strict comparison
of divided-power algebras, a common admissible source, or a
contracting homotopy in their localization. It rules out a
specified family of proposed \emph{linear} invariants assembled
from infinitely many coefficient primes.

\paragraph{ARCH-136-LINEAR-PRIME-TAIL.}
Fix a finite nonempty family $\mathscr C$ of admissible integral
theories, and integers $a,d\geq0$. There is $W=W(\mathscr C,a,d)$
with the following property. For any $C\in\mathscr C$, any
simplicial set $X$, any weight $w\geq W$, and any bounded complex
$P$ of flat abelian groups concentrated in degrees $[-a,b]$
(the finite upper bound $b$ need not be fixed), the actual natural
comparison zigzag induces isomorphisms through degree $d$ between
\[
 C^{w,*}(X)\otimes^{\mathbf L}_{\mathbb Z}P
 \quad\text{and}\quad
 \mathcal E_X(\mathbb Z)\otimes^{\mathbf L}_{\mathbb Z}P
                         \otimes_{\mathbb Z}\mathbb ZT^{[w]},
 \qquad
 \mathcal E_X(\mathbb Z)=
 R\operatorname{Hom}_{s\operatorname{Ab}}
       (\mathbb Z[X],\underline{\mathbb Z}).                 \tag{136.LP1}
\]
The bound is independent of $X$, $P$, its entries, every integer
appearing in its differentials, and the weight $w\geq W$.
These comparisons are natural for strict arrows of theories,
maps of simplicial sets, maps of coefficient complexes, and
coefficient multiplication maps between the specified weights.

Consequently, choose any family of such data indexed by a set,
with all weights at least $W$ and with the same lower amplitude
bound $-a$. Taking their products preserves the comparisons
through degree $d$. Taking a countable homotopy inverse limit
also preserves them, whenever the data form a tower and the
transition maps commute with the actual comparison zigzags.
Thus a product, a prime-diagonal subsequence, or a tower
$\lim^1$ class obtained entirely from these natural bounded
derived coefficient comparisons is theory-independent in this
range. In particular the primes or prime-power exponents may
be unbounded.

\begin{proof}
For an admissible $C$, put $M=C^{w,*}(X)$ and
$E=\mathcal E_X(\mathbb Z)\otimes\mathbb ZT^{[w]}$.
AQ-136-COMPARISON supplies the actual natural arrows
\[
 M\longrightarrow
 R\Gamma_X(C^{w,*})\longleftarrow E,
 \qquad
 \Gamma_X(F)=\operatorname{Hom}_{s\operatorname{Ab}}
                       (\mathbb Z[X],F).                 \tag{136.LP2}
\]
For one bound independent of $X$, each arrow is an isomorphism
on integral cohomology through degree $d+a$. Take the maximum
of these bounds for the finite family $\mathscr C$.
All complexes in (136.LP2) have zero cohomology in negative
degrees. The right arrow is induced by the coefficient map;
the left arrow is the ordinary-to-derived-section map. This
explains all the stated naturality, including coefficient
multiplication $T^{[v]}:C^{w,*}\to C^{w+v,*}$.

Here is the precise bounded-amplitude argument, without an
unjustified tensoring of cohomology isomorphisms. For any complex
$K$ and $D=d+a$, use the canonical triangle
\[
 \tau_{\leq D}K\longrightarrow K\longrightarrow
                          \tau_{\geq D+1}K\longrightarrow.
\]
A bounded complex of flat modules is $K$-flat. Moreover,
tensoring a complex in $D^{\geq D+1}(\mathbb Z)$ with $P$
places it in $D^{\geq D+1-a}(\mathbb Z)$: filter the finite
complex $P$ by its terms; each flat term preserves the lower
cohomological bound, shifted by its degree, and extensions
preserve that bound. It follows that
\[
 H^j((\tau_{\leq D}K)\otimes^{\mathbf L}P)
       \longrightarrow H^j(K\otimes^{\mathbf L}P)
             \quad\text{is an isomorphism for }j\leq d.
\]
An arrow inducing isomorphisms through $D$ induces an
isomorphism on $\tau_{\leq D}$; apply the displayed observation
to each arrow of (136.LP2). This proves (136.LP1), with the
claimed single bound. For $P=[\mathbb Z\xrightarrow{n}\mathbb Z]$
in degrees $-1,0$, this recovers the uniform-in-$n$ statement
of AQ-136-DERIVED-COMPATIBILITY. It is the same actual
comparison, not an independently chosen family of abstract
cohomology identifications.

Products are exact in abelian groups. Thus the cohomology of
a product of complexes is the product of the cohomologies,
and taking products of the two comparison arrows preserves
their isomorphisms through degree $d$.

For a countable tower $K_0\leftarrow K_1\leftarrow\cdots$,
its homotopy inverse limit is the fiber of
\[
 1-s:\prod_{i\geq0}K_i\longrightarrow\prod_{i\geq0}K_i,
\]
where $s$ is induced by the transition maps. The fiber long
exact sequence gives, in each degree $j$, the natural exact
sequence
\[
 \begin{split}
 0&\to\operatorname{coker}(1-s:\prod_iH^{j-1}K_i
                              \to\prod_iH^{j-1}K_i)\\
  &\to H^j(\operatorname{holim}_i K_i)
 \to\ker(1-s:\prod_iH^jK_i\to\prod_iH^jK_i)\to0.
 \end{split}
                                                               \tag{136.LP3}
\]
The two outer terms are the usual $\lim^1 H^{j-1}K_i$ and
$\lim H^jK_i$. An objectwise comparison isomorphism in degrees
$j-1,j$ identifies both of them. The short five lemma
therefore proves the assertion for $j\leq d$. Negative degrees
cause no problem: all the relevant complexes have the common
lower cohomological bound $-a$, and the preceding tensor
argument applies to every degree at most $d$.
\end{proof}

\paragraph{What this does and does not exclude.}
The inherited fixed-prime Bockstein obstruction
AQ-136-HIGHER-COMPLETION concerns maps out of a specified
\emph{nonadmissible} partial source. It remains valid.
The present lemma says that merely letting that prime escape
to infinity, and packaging the linear derived coefficient
classes by products or a tower, cannot distinguish two
admissible theories: the comparison bound is chosen before
all those primes. A proposed counterexample would need
additional nonlinear defining-system data or operations
outside this bounded derived coefficient construction.
No assertion is made for unbounded negative amplitude, for
ordinary reduction of nonflat forms, or for transition maps
which fail to commute with the actual comparison.
